\documentclass[sigconf,nonacm,11pt]{acmart}

\usepackage{graphicx}
\usepackage{amsmath}

\setcopyright{none}
\settopmatter{printacmref=false, printccs=false, printfolios=true}
\acmConference[Working Paper]{Solana MEV Study}{2026}{Working Paper}

\begin{document}
\title{Identification and Anatomy of Sandwich Attacks on Solana Jito}
\author{Anonymous}
\affiliation{\institution{Working Paper}\city{}\country{}}
\maketitle

\section{Sandwich-Attack Identification}
\label{sec:detection}

Our dataset is the join of Jito bundle metadata and Helius enriched
transaction data, comprising $3{,}783{,}142$ transactions across
$1{,}260{,}686$ bundles. Each transaction supplies wallet-level flow
through \texttt{tokenTransfers} and \texttt{nativeTransfers}.

For every bundle of length $\geq 3$, we enumerate ordered triples
$(i, k, j)$ and construct the wallet--asset signed-flow vector
$\Delta_x(w, m) = \mathrm{inflow}_x(w, m) - \mathrm{outflow}_x(w, m)$
per transaction.
A wallet $W$ is flagged as a candidate attacker if there exists an
asset pair $(b, q)$ such that $W$ executes mutually reversed swaps on
$(b, q)$ between $x_i$ and $x_j$ (i.e., $\Delta_{x_i}(W, b)$ and
$\Delta_{x_j}(W, b)$ have opposite signs and the quote flips
correspondingly).
We deliberately do \emph{not} require both legs to share a
\texttt{feePayer}: in our sample $99.7\%$ of attackers use a dedicated
tipper wallet, and any signer-based pre-filter would discard the
overwhelming majority of true positives.

Each candidate $(W, i, k, j)$ must additionally pass three economic
filters: (i) \textbf{symmetry}, $|Q_f - Q_b|/\max(|Q_f|, |Q_b|) \leq 0.10$,
ruling out one-sided liquidations and inventory dumps;
(ii) \textbf{victim direction}: at least one wallet $V \neq W$ in some
intermediate $x_k$ has a quote-asset flow aligned with $W$'s front leg;
(iii) \textbf{positive base profit}.
Two cleaning rules remove residual artefacts:
R1 drops bundles whose flagged attacker is also a victim signer or
aligned wallet (self-attacks);
R2 blacklists wallets that appear as candidate attacker in $\geq 200$
bundles with balanced direction ($\min(b,s)/\max(b,s) \geq 0.30$) and
mean bundle profit within $\pm 0.01$ SOL---the empirical signature of
two-sided market makers rather than opportunistic attackers.

The pipeline identifies
\textbf{$11{,}591$ confirmed Pattern-A sandwich attacks},
spanning $12{,}916$ unique slots and $161$ unique attacker wallets;
$5{,}060$ ($38.4\%$) realise positive net SOL profit.
This count is approximately \emph{two orders of magnitude} larger
than previously reported Pattern-A volumes on Solana, the gap being
attributable almost entirely to dropping the implicit
\texttt{feePayer}-equality assumption.

\begin{figure}[t]
  \centering
  \includegraphics[width=\columnwidth]{figures/sandwich_anatomy_col_en.pdf}
  \caption{Four-axis anatomy of $11{,}591$ confirmed sandwiches.}
  \label{fig:anatomy}
\end{figure}

Figure~\ref{fig:anatomy} dissects the confirmed sample along four
dimensions.
Panel~(A) is the joint decomposition of price-edge versus
net-SOL profitability: although $75\%$ of attackers correctly position
on the price spread, only $36\%$ realise positive net SOL profit; the
doubly-profitable cell accounts for just $14.5\%$ of attempts.
Most ``price-edge wins'' are erased by pool slippage and residual
base-token inventory.
Panel~(B) projects both $|\textit{profit}|$ and $\textit{tip}+\textit{fee}$
onto the capital axis on log--log scales: the profit slope is
$\approx 0.99$ (nearly linear in capital), whereas the cost slope is
$\approx -0.28$ (essentially flat).
Marginal cost is close to zero, granting capital-rich attackers a
structural scale advantage.
Panel~(C) shows that the market is highly concentrated
(Gini $= 0.78$, top $1\%$ executes $11\%$ and top $10\%$ executes
$79\%$), while two insets reveal that $86\%$ of attacks are
front-buyer attacks and $99.2\%$ of bundles target only one victim.
Panel~(D) plots the joint density of execution symmetry
(\texttt{sym\_ratio}) and SOL profit:
$sym\_ratio < 0.02$ yields positive expected profit, whereas
$sym\_ratio > 0.05$ is loss-making on average---in the absence of
bidding wars, what separates winners from losers is execution
precision rather than tip aggression.

\end{document}
